\subsubsection{Ramanujan--视界对称框架下的零点“视界纯度”判据与一条可执行的谱约化}\label{subsubsec:xi-ramanujan-horizon-purity-spectral-reduction}
本段在既有的两条接口之上抽取三条互补的结构结论，并将其组织为一条可直接引用的闭环：其一把 RH 完全约化为单标量 Jensen 缺陷在视界极限的消失；其二把有限维 self-inversive 完成化多项式的“单位圆根/Ramanujan”条件译码为相同的缺陷消失判据；其三以复分析意义下的 Hurwitz 零点定理把“构造一列视界纯的逼近对象”约化为推出 RH 的充分条件。上述结论均在文稿既有编号链内闭合：圆盘化完成化与 Jensen 缺陷见附录 \ref{app:unit-circle-phase-gate} 的 \S\ref{subsec:unit-disk-jensen-defect}，Ramanujan 相位化原则见 \S\ref{subsubsec:ramanujan-phase-lift-symmetry}，而固定切片承载与 Blaschke/点谱缺陷见 \S\ref{subsubsec:selfdual-blaschke-defect-fixed-slice}。
\par
为避免歧义，本文在第 \ref{sec:principles} 节使用的 Hurwitz/Markov 逼近口径与本段所用的 Hurwitz \emph{零点定理}（复分析）属于不同语境；本段的“Hurwitz”始终指后者。

\paragraph{记号与视界紧化}
沿用附录 \ref{app:unit-circle-phase-gate} 的圆盘化完成化接口：Cayley 紧化 \(\mathcal{C}:\mathbb{H}\to\mathbb{D}\) 及其逆映射由 \eqref{eq:app-cayley-disk} 给出，
并令
\begin{equation}\label{eq:xi-ramanujan-horizon-s-of-w}
s(w):=\frac12+\mathrm{i}\,\mathcal{C}^{-1}(w)\qquad(w\in\mathbb{D}),
\end{equation}
则圆盘化完成化函数
\begin{equation}\label{eq:xi-ramanujan-horizon-Fzeta}
F_{\zeta}(w):=\Xi_{\zeta}\!\bigl(s(w)\bigr)\qquad(w\in\mathbb{D})
\end{equation}
在 \(\mathbb{D}\) 内解析（附录式 \eqref{eq:app-Fzeta-disk}）。
对 \(0<\varrho<1\)，Jensen 缺陷（定义 \ref{def:app-jensen-defect}）为
\begin{equation}\label{eq:xi-ramanujan-horizon-jensen-defect}
D_{\zeta}(\varrho)
:=\frac{1}{2\pi}\int_{0}^{2\pi}\log\bigl|F_{\zeta}(\varrho e^{\mathrm{i}\theta})\bigr|\,\dd\theta-\log|F_{\zeta}(0)|\ \in\RR_{\ge 0}.
\end{equation}
在本文术语中，\(D_{\zeta}(\varrho)\) 可被理解为视界层 \(|w|=\varrho\) 上的“视界熵缺陷”；该术语仅用于对齐解释性语言，与证明无关。

\begin{theorem}[RH 的单标量“视界纯度”判据]\label{thm:xi-horizon-purity-criterion}
以下命题等价：
\begin{enumerate}
  \item[(i)] RH 成立；
  \item[(ii)] \(\displaystyle \lim_{\varrho\uparrow 1}D_{\zeta}(\varrho)=0\)；
  \item[(iii)] \(\displaystyle \limsup_{\varrho\uparrow 1}D_{\zeta}(\varrho)=0\)。
\end{enumerate}
\end{theorem}

\begin{proof}
\textbf{(i)\(\Rightarrow\)(ii).}
由推论 \ref{cor:app-rh-iff-jensen-defect-zero}，RH 成立当且仅当 \(D_{\zeta}(\varrho)\equiv 0\) 于 \((0,1)\)，从而 \(\lim_{\varrho\uparrow 1}D_{\zeta}(\varrho)=0\)。
\par
\textbf{(ii)\(\Rightarrow\)(iii).}
显然。
\par
\textbf{(iii)\(\Rightarrow\)(i).}
反设 RH 不成立。由定理 \ref{thm:app-rh-iff-disk-zero-free}，存在盘内零点 \(a\in\mathbb{D}\) 使 \(F_{\zeta}(a)=0\)。
则由命题 \ref{prop:app-jensen-single-zero-lower-bound}（单零点 Jensen 下界）得
\[
\liminf_{\varrho\uparrow 1}D_{\zeta}(\varrho)\ \ge\ \log(1/|a|)\ >\ 0,
\]
与 \(\limsup_{\varrho\uparrow 1}D_{\zeta}(\varrho)=0\) 矛盾。
故盘内零点不存在，即 \(F_{\zeta}\) 在 \(\mathbb{D}\) 内无零点；再由定理 \ref{thm:app-rh-iff-disk-zero-free} 得 RH 成立。证毕。
\end{proof}

\begin{remark}[与可审计证书链的对齐]\label{rem:xi-horizon-purity-audit-alignment}
定理 \ref{thm:xi-horizon-purity-criterion} 的内容与附录定理 \ref{thm:app-rh-iff-jensen-defect-vanishing-sequence} 等价；其用途在于：将“盘内零点事件”完全压缩为单参标量 \(D_{\zeta}(\varrho)\) 的边界极限行为，并与正文中“缺陷上界序列 \(\Rightarrow\) 排斥半径趋一 \(\Rightarrow\) RH”的证书链（定理 \ref{thm:dephys-defect-repulsion-radius}）在形式上直接同构。
\end{remark}

\paragraph{Ramanujan 谱界的视界熵形式：self-inversive 多项式的 Jensen 判据}
为把有限维 Ramanujan/RH-sharp 的“单位圆根”条件与缺陷判据对齐，引入多项式层面的 Jensen 缺陷。

\begin{definition}[多项式的 Jensen 缺陷]\label{def:xi-poly-jensen-defect}
令 \(P(y)\in\CC[y]\) 为非零多项式且 \(P(0)\neq 0\)。
对 \(0<\varrho<1\) 定义
\[
D_{P}(\varrho)
:=\frac{1}{2\pi}\int_{0}^{2\pi}\log\bigl|P(\varrho e^{\mathrm{i}\theta})\bigr|\,\dd\theta-\log|P(0)|.
\]
\end{definition}

\begin{theorem}[self-inversive 完成化多项式的“视界熵”判据]\label{thm:xi-selfinversive-jensen-criterion}
设 \(P(y)\in\CC[y]\) 为 self-inversive 多项式：存在 \(\phi\in\RR\) 与 \(d=\deg P\) 使
\begin{equation}\label{eq:xi-selfinversive}
y^{d}\overline{P(1/\overline{y})}=e^{\mathrm{i}\phi}P(y).
\end{equation}
并设 \(P(0)\neq 0\)。
令 \(D_P(\varrho)\) 如定义 \ref{def:xi-poly-jensen-defect}。
则以下断言等价：
\begin{enumerate}
  \item[(i)] \(P\) 的全部零点都满足 \(|y|=1\)；
  \item[(ii)] \(D_{P}(\varrho)\equiv 0\) 对一切 \(0<\varrho<1\) 成立；
  \item[(iii)] \(\displaystyle \lim_{\varrho\uparrow 1}D_{P}(\varrho)=0\)。
\end{enumerate}
\end{theorem}

\begin{proof}
\textbf{(i)\(\Rightarrow\)(ii).}
若 \(P\) 的零点皆在 \(|y|=1\)，则对任意 \(\varrho<1\)，圆盘 \(|y|<\varrho\) 内无零点。
由命题 \ref{prop:app-jensen-formula-defect} 的 Jensen 求和式（对解析函数 \(P\) 在 \(|y|<\varrho\) 的应用）得 \(D_P(\varrho)=0\)。
\par
\textbf{(ii)\(\Rightarrow\)(iii).} 显然。
\par
\textbf{(iii)\(\Rightarrow\)(i).}
若存在零点 \(a\) 满足 \(|a|<1\)，则由命题 \ref{prop:app-jensen-single-zero-lower-bound}（应用于 \(F=P\)）得对任意 \(\varrho\in(|a|,1)\) 有
\(
D_P(\varrho)\ge \log(\varrho/|a|)
\)，从而
\(
\liminf_{\varrho\uparrow 1}D_P(\varrho)\ge \log(1/|a|)>0
\)，与 (iii) 矛盾；故 \(P\) 在 \(|y|<1\) 内无零点。
\par
由 self-inversive 条件 \eqref{eq:xi-selfinversive}，若 \(b\) 为零点且 \(|b|>1\)，则 \(1/\overline{b}\) 亦为零点且满足 \(|1/\overline{b}|<1\)，与“盘内无零点”矛盾；故亦无 \(|y|>1\) 的零点。
因此 \(P\) 的全部零点必满足 \(|y|=1\)。证毕。
\end{proof}

\begin{corollary}[有限维 toy-RH/Ramanujan 的缺陷译码]\label{cor:xi-finite-ramanujan-defect-translation}
在 \S\ref{subsec:xi-completion-audit} 的完成化口径中，完成化多项式 \(\Phi_L\) 为 self-inversive（命题 \ref{prop:xi-rh-cohn-derivative}）。
因此，\(\Xi\)-RH 的“单位圆根判据”（推论 \ref{cor:xi-rh-unit-circle-reciprocal}）可等价重写为 \(\Phi_L\) 的 Jensen 缺陷 \(D_{\Phi_L}(\varrho)\equiv 0\)（或 \(\lim_{\varrho\uparrow 1}D_{\Phi_L}(\varrho)=0\)）。
该等价与 \S\ref{subsubsec:ramanujan-phase-lift-symmetry} 的“Ramanujan 谱界 \(\Leftrightarrow\) 相位化闭包”同型：两者都把谱侧约束转译为“无径向/幅度泄漏”的可核验判据，而非坐标轴置换。
\end{corollary}

\paragraph{谱逼近到 RH：Hurwitz 零点定理的约化}
下述约化把“证明 RH”转写为一个可操作的构造性问题：在圆盘完成化口径下，构造一列在 \(\mathbb{D}\) 内无零点的逼近对象并证明其在紧集上一致收敛到 \(F_{\zeta}\)。

\begin{theorem}[Ramanujan--视界逼近原理：Hurwitz 零点定理约化]\label{thm:xi-hurwitz-horizon-reduction}
设 \(\{F_N\}_{N\ge 1}\) 为一列解析函数 \(F_N:\mathbb{D}\to\CC\)，并满足：
\begin{enumerate}
  \item \textbf{（视界纯度）}对每个 \(N\)，\(F_N\) 在 \(\mathbb{D}\) 内无零点；
  \item \textbf{（自对偶完成化一致性）}每个 \(F_N\) 均来自某一自对偶完成化行列式/核对象（从而其对偶对合与边界实性与本文的 unitary slice 口径兼容，见 \S\ref{subsec:xi-completion-audit}）；
  \item \textbf{（局部一致收敛）}\(F_N\to F_{\zeta}\) 在 \(\mathbb{D}\) 的每个紧子集上一致收敛，其中 \(F_{\zeta}\) 由 \eqref{eq:xi-ramanujan-horizon-Fzeta} 定义。
\end{enumerate}
则 RH 成立。
\end{theorem}

\begin{proof}
由 Hurwitz 零点定理（复分析；见 \cite[Chapter~4]{Ahlfors1979ComplexAnalysis}）：若解析函数列在域内局部一致收敛且每个函数在该域内无零点，则极限函数要么恒为零，要么亦无零点。
由附录 \ref{app:unit-circle-phase-gate}，\(F_{\zeta}(0)=\Xi_{\zeta}(-\tfrac12)\neq 0\)，故 \(F_{\zeta}\not\equiv 0\)。
于是 \(F_{\zeta}\) 在 \(\mathbb{D}\) 内无零点。
再由定理 \ref{thm:app-rh-iff-disk-zero-free} 得 RH 成立。证毕。
\end{proof}

\paragraph{闭环式总结（结构等价而非修辞）}
定理 \ref{thm:xi-horizon-purity-criterion} 把 RH 压缩为 Jensen 缺陷的单标量极限判据；
定理 \ref{thm:xi-selfinversive-jensen-criterion} 给出有限维 self-inversive 完成化对象的同型缺陷判据，并在推论 \ref{cor:xi-finite-ramanujan-defect-translation} 中与 toy-RH/Ramanujan 的单位圆根判据对齐；
定理 \ref{thm:xi-hurwitz-horizon-reduction} 则把“构造视界纯的 Ramanujan 逼近”升级为推出 RH 的充分条件。
在本文的协议语义语言下，这三条结论共同给出一个可审计闭环：
\[
\text{自对偶（固定切片）}\ +\ \text{Ramanujan 相位化（谱界）}
\ \Longrightarrow\ 
\text{视界纯度（缺陷消失）}
\ \Longrightarrow\ 
\text{盘内无零点}
\ \Longrightarrow\ 
\text{RH}.
\]
\endinput
